% SPDX-License-Identifier: Apache-2.0
% covers: Plic1, Plic2, Plic3, Plic4, Plic5, Plic6, Plic7, Plic8
\datasheet{Plic}{Platform-level interrupt controllers: one to eight sources and one target, at the RISC-V PLIC's offsets.}

\dsfacts{\code{lib/parts/src/plic.rs}; generator \code{plic\_text} in \code{lib/macros/lib.rs}}%
{\code{txhdl\_parts::plic::Plic1} to \code{Plic8}}%
{\code{//docs:vreteno}}

\subsection*{Function}

A controller gathers the interrupt lines of several sources into one
line for one target, a hart's machine mode, and lets software ask which
source to serve. Each source \textit{i} is an input line,
\code{src}\textit{i}, numbered from 1. The output \code{irq} is the
target's line. Software reaches the controller as an AXI-Lite
peripheral: reads come on \code{ar} and are answered on \code{r}, and
writes come on \code{aw} and \code{w} and are answered on \code{b}. It
sits behind a \code{LiteBridge1}, which sends it only its own range.
The input \code{rst} clears the requests, the line, the priorities, the
enable bits and the threshold.

\code{plic!(PlicN, N)} writes one unit per count, since the lowering
reads a body and not a loop over ports. \code{plic.rs} writes
\code{Plic1} to \code{Plic8}; the macro takes counts from 1 to 31, the
most one word of pending bits holds. \code{plic.rs} also names the
offsets, \code{PENDING}, \code{ENABLE}, \code{THRESHOLD} and
\code{CLAIM}, and \code{priority(i)}.

\subsection*{Parameters}

\begin{center}\footnotesize
\begin{tabular}{@{}lp{0.7\textwidth}@{}}
\toprule
Parameter & Meaning \\
\midrule
\textit{N} & The count of sources, in the type's name. \\
\code{EDGE} & A bit per source: bit \textit{i} set makes source
\textit{i} ask on a rising edge rather than while its line is high.
Bit 0 and bits above \textit{N} are ignored. The tables use 0, as the
board does. \\
\bottomrule
\end{tabular}
\end{center}

\subsection*{Register map}

Offsets from the controller's base, which are the RISC-V PLIC's for
one target. The controller decodes the low 22 bits of an address. On
the board the base is \code{0x0c00\_0000}.

\begin{center}\footnotesize
\begin{tabular}{@{}llp{0.55\textwidth}@{}}
\toprule
Offset & Word & Access \\
\midrule
\code{0x4} $\times$ \textit{i} & Source \textit{i}'s priority, 3 bits &
Read and write, for \textit{i} from 1 to \textit{N}. \\
\code{0x1000} & The pending bits, bit \textit{i} for source \textit{i} &
Read only. \\
\code{0x2000} & The enable bits, the same way & Read and write; bit 0
does not stick. \\
\code{0x20\_0000} & The threshold, 3 bits & Read and write. \\
\code{0x20\_0004} & Claim and complete & A read claims; a write
completes. \\
\bottomrule
\end{tabular}
\end{center}

Every other offset, offset 0 among them, reads as zero and ignores a
write.

\dstables{Plic}

\subsection*{Behaviour}

\begin{itemize}
\item \textbf{The gateway.} A level source asks while its line is high.
An edge source asks when its line is high and was low a cycle before
(\code{prev}), or when it holds an edge (\code{held}).
\item A request goes forward when its source is not \code{active}.
Going forward sets the source's bits in \code{pending} and in
\code{active}.
\item An edge that comes while its source is \code{active} sets its bit
in \code{held}, one deep, and it asks as soon as the source is no
longer active.
\item \textbf{The choice.} A source is a candidate when it is pending,
enabled, and its priority is above \code{threshold}. The best
candidate has the highest priority, and the lowest number wins a tie.
A priority of 0 is never above a threshold, so it never interrupts.
\item \textbf{A claim.} A read of \code{0x20\_0004} answers with the
best candidate's number, or 0 for none, and clears that source's
\code{pending} bit. The source stays \code{active}.
\item \textbf{A complete.} A write of a number from 0 to \textit{N} to
\code{0x20\_0004} clears that source's \code{active} bit. A larger
number changes nothing. A level source whose line is still high then
asks again at once.
\item \textbf{The line.} \code{asserted} is set each cycle when there
is a best candidate, and \code{irq} follows \code{asserted}, so the
line is a register's output and lags the choice by a cycle.
\item \textbf{The bus.} A read is answered in the cycle it is taken, and
a write is taken when its address and its word are both there, and
answered \code{Okay} at once. Every answer is \code{Okay}.
\end{itemize}

\subsection*{Verification}

\begin{itemize}
\item In \code{plic}, in \code{//lib/parts:parts\_test}, against a
\code{Plic3}:
\begin{itemize}
\item \code{the\_highest\_priority\_is\_claimed\_first\_and\_the\_lowest\_number\_wins\_a\_tie};
\item \code{a\_priority\_at\_or\_below\_the\_threshold\_does\_not\_interrupt};
\item \code{a\_claimed\_source\_asks\_again\_only\_after\_its\_complete};
\item \code{an\_edge\_source\_asks\_once\_per\_edge\_and\_a\_level\_source\_while\_high}.
\end{itemize}
\code{the\_eight\_source\_controller\_lowers} writes \code{Plic8}'s
Verilog.
\item \code{ex\_plic} runs \code{Plic3<0b100>} against a host that
claims and completes, and the build simulates its netlist against that
run under nvc and Verilator: \code{//docs:plic\_sim\_plic3\_tb\_test}
and \code{//docs:plic\_vsim\_test}.
\item \code{input\_comes\_by\_interrupt\_one\_byte\_each}, in
\code{//cpu/vreteno:board\_test}, runs the board's \code{Plic2} under a
compiled program that takes the serial port's input by interrupt.
\end{itemize}

\subsection*{Used by}

\code{Plic2<0>} in the board's design, \code{Board}, as its field
\code{plic} behind the bridge \code{pplic}: source 1 is the serial
port's receive interrupt, source 2 the board's \code{irq} input, and
\code{irq} is the core's external interrupt. \code{Plic3<0b100>} in the
example \code{ex\_plic}.

\subsection*{Limits}

It has one target and no per-target context beyond it. Priorities are 3
bits, so from 0 to 7. A write ignores its strobe and writes the whole
field. The controller checks no address above bit 21, and relies on its
bridge for the rest. An edge source holds one edge during a service; a
second is lost.
